\chapter*{Введение}
\addcontentsline{toc}{chapter}{Введение}

\textbf{Актуальность темы.} Моделирование распространения звуковых волн в одномерных волноводах переменного сечения представляет собой классическую задачу технической акустики, которая находит широкое применение во многих прикладных областях. К ним относятся акустика музыкаческих инструментов (в частности, духовых рупоров), проектирование акустических систем и волноводных концентраторов, а также моделирование речевого тракта человека для задач синтеза речи. Одномерное приближение на основе уравнения Вебстера позволяет существенно сократить вычислительные затраты по сравнению с полномерным трехмерным моделированием уравнений Навье--Стокса или Гельмгольца при сохранении высокой физической точности в области низких и средних частот.

Существующие подходы к численному расчету одномерного акустического поля делятся на методы матриц передачи (TMM) для кусочно-постоянной или кусочно-конической аппроксимации геометрии, метод матриц длинных линий (Transmission Line Matrix, TLM), а также авторегрессионные схемы (ARMA), позволяющие свести расчет к быстрым полиномиальным операциям. Каждый из этих подходов обладает уникальным балансом вычислительной сложности, точности и возможностей расширения. В связи с этим актуальной задачей является проведение систематического исследования и сравнительного анализа данных методов на репрезентативном наборе тестовых геометрий.

\textbf{Цель работы} --- разработка кроссплатформенного программного комплекса для численного моделирования акустических полей в волноводах переменного сечения и проведение комплексного сравнительного анализа методов TMM (цилиндрического и конического), авторегрессионного моделирования ARMA и метода TLM по критериям точности, сходимости и времени выполнения.

Для достижения поставленной цели необходимо решить следующие \textbf{задачи}:
\begin{enumerate}
    \item Изучить теоретические основы одномерной волновой акустики и математический аппарат уравнения Вебстера.
    \item Реализовать методы расчета передаточных матриц для цилиндрических и конических сегментов (TMM).
    \item Разработать эффективный решатель на основе авторегрессионного скользящего среднего (ARMA) для ускорения частотных расчетов.
    \item Разработать эффективный решатель на основе метода матриц длинных линий (TLM) на языке C++ с возможностью интеграции в Python-окружение.
    \item Провести серию численных экспериментов по верификации реализованных решателей на регулярных и сложных геометриях волноводов.
    \item Выполнить комплексный сравнительный анализ вычислительной эффективности и точности реализованных методов.
\end{enumerate}

\textbf{Научная новизна} работы заключается в систематическом сравнении классических частотных методов TMM с авторегрессионным подходом ARMA в z-области и временным методом TLM для широкого класса гладких, кусочно-непрерывных и резонансных геометрий волноводов.


